% C4 Architecture Visualization Documentation Suite | Original educational content
% Standalone LaTeX root: no external style file, no network, no shell escape.
\documentclass[10pt,letterpaper]{article}
\usepackage[T1]{fontenc}
\usepackage[utf8]{inputenc}
\usepackage{lmodern,microtype}
\usepackage[letterpaper,margin=0.82in,headheight=15pt,headsep=18pt,footskip=26pt]{geometry}
\usepackage[table]{xcolor}
\usepackage{booktabs,tabularx,array,longtable}
\usepackage{enumitem}
\usepackage{fancyhdr,lastpage}
\usepackage{titlesec}
\usepackage{hyperref,xurl}
\usepackage{needspace}
\usepackage{listings}
\usepackage{amssymb}
\definecolor{C4Navy}{HTML}{17364B}
\definecolor{C4Teal}{HTML}{137C85}
\definecolor{C4Gray}{HTML}{526775}
\definecolor{C4Light}{HTML}{EEF4F5}
\definecolor{C4Line}{HTML}{CAD8DE}
\color{C4Navy}
\hypersetup{colorlinks=true,linkcolor=C4Teal,urlcolor=C4Teal,citecolor=C4Teal}
\setlength{\parindent}{0pt}
\setlength{\parskip}{5pt}
\setlength{\emergencystretch}{2em}
\raggedbottom
\setlist[itemize]{leftmargin=1.4em,topsep=3pt,itemsep=3pt}
\setlist[enumerate]{leftmargin=1.6em,topsep=3pt,itemsep=3pt}
\titleformat{\section}{\large\bfseries\sffamily\color{C4Navy}}{\thesection}{0.6em}{}
\titleformat{\subsection}{\normalsize\bfseries\sffamily\color{C4Teal}}{\thesubsection}{0.6em}{}
\titlespacing*{\section}{0pt}{16pt}{6pt}
\titlespacing*{\subsection}{0pt}{9pt}{4pt}
\pagestyle{fancy}
\fancyhf{}
\fancyhead[L]{\scriptsize\sffamily C4 / ARCHITECTURE VISUALIZATION}
\fancyhead[R]{\scriptsize\sffamily REFERENCE / EDUCATION}
\fancyfoot[L]{\scriptsize\color{C4Gray}ORIGINAL ENGINEERING SUITE / 2026}
\fancyfoot[R]{\scriptsize\thepage\ / \pageref*{LastPage}}
\renewcommand{\headrulewidth}{0.4pt}
\newcolumntype{Y}{>{\raggedright\arraybackslash}X}
\newcommand{\RuleLine}{\par\noindent\color{C4Line}\rule{\linewidth}{0.45pt}\color{C4Navy}\par}
\newcommand{\WritingSpace}[1]{\par\vspace{2pt}\noindent\fcolorbox{C4Line}{white}{\parbox[t][#1][t]{0.94\linewidth}{\hspace{1pt}}}\par\vspace{7pt}}
\newcommand{\Field}[1]{\par\textbf{#1}\par\WritingSpace{0.57in}}
\newcommand{\Prompter}[1]{\par{\small\color{C4Teal}\textbf{Authoring guidance:} #1}\par}
\newcommand{\SourceNote}[1]{\par\smallskip\noindent\colorbox{C4Light}{\parbox{0.94\linewidth}{\footnotesize\textbf{PROVENANCE.} #1}}\par\medskip}
\lstset{basicstyle=\ttfamily\footnotesize,breaklines=true,columns=fullflexible,showstringspaces=false,frame=single,rulecolor=\color{C4Line}}
\newcommand{\DocStart}[3]{%
\begin{document}\begin{titlepage}
\vspace*{0.55in}
{\sffamily\bfseries\small\color{C4Teal}#1}\par\vspace{0.24in}
{\sffamily\bfseries\fontsize{26}{31}\selectfont\color{C4Navy}#2}\par\vspace{0.18in}
{\large\color{C4Gray}#3}\par\vspace{0.4in}
\textcolor{C4Teal}{\rule{\linewidth}{1.15pt}}\par\vspace{0.24in}
\textbf{Document control}\par
Edition 1.0\quad|\quad 30 October 2026\quad|\quad Original educational engineering material\par
\vspace{0.17in}
\textbf{Source boundary}\par
Simon Brown, \emph{The C4 Model: Visualizing Software Architecture} (O'Reilly, 2026), is the primary study source. The suite's exercises, forms, crosswalks, process rules, sample systems, and scoring models are independently authored supplements; they are not represented as the author's prescriptions.\par
\vfill
{\small\color{C4Gray}Not an ISO/IEC/IEEE 42010 conformance certificate. Not a substitute for the published standard or for review of deployed systems.}\par
\end{titlepage}
}
\DocStart{PREREQUISITE CURRICULUM}{C4 Architecture Visualization Curriculum}{Ten-week, outcome-based learning path}
\tableofcontents\clearpage
\SourceNote{Reading weeks follow Brown (2026); sequencing, assessment and lab schedule are suite-designed.}
\section{Entry diagnostics and prerequisites}
\begin{itemize}
\item Software architecture fundamentals: responsibilities, interfaces, static vs. runtime structure.
\item Git source-control basics and clear technical writing.
\item A small project codebase or the supplied fictional Civic Permit Portal scenario.
\item Working PDF/LaTeX compilation and repository PlantUML renderer for diagram labs.
\end{itemize}

\section{Weekly crosswalk}
\begin{tabularx}{\linewidth}{@{}YYY@{}}\toprule
\textbf{Week / primary topic} & \textbf{Reading} & \textbf{Evidence / checkpoint} \\ \midrule
1: Architecture communication & Book Chapter 1 & Weak diagram critique; Find 5 ambiguity issues \\
2: C4 abstraction model & Book Chapter 2 & Element classification matrix; Distinguish container and deployment node \\
3: Boundaries and context & Book Chapter 3 & Context diagram; Lab 01–02; Stakeholder peer review \\
4: Container decomposition & Book Chapter 4 & Container view; Lab 03; Map context relationships \\
5: Component and code zoom & Book Chapters 5–6 & Component mapping; Lab 04; Map component -> code \\
6: Runtime behavior & Book Chapter 7 & Dynamic happy/failure paths; Lab 05; Numbered, evidenced interactions \\
7: Runtime placement & Book Chapter 8 & Dev and prod deployment; Lab 06; Reconcile logical to deployed \\
8: Enterprise scale & Book Chapter 9 & Landscape; Lab 07; Organization scope verified \\
9: Visual precision and scaling & Book Chapters 10–11 & Notation audit; async architecture; Lab 08; Cross-diagram consistency \\
10: Operating the documentation system & Book Chapter 12 and Appendix & PR review, scorecards; Lab 09; Accepted capstone evidence package \\
\bottomrule\end{tabularx}


\section{Dependency graph}
Communication -> common abstraction vocabulary -> context -> container -> component -> code. Dynamic views depend on stable element names; deployment depends on the logical container inventory. Landscape views aggregate per-system contracts. Notation and governance cut across the whole path.


\section{Capstone and grading}
\begin{itemize}
\item Deliver a set of minimum useful diagrams with source versions and a purpose statement for each.
\item Show two traceability paths from stakeholder concerns to code or deployment evidence.
\item Complete one architecture-diagram scorecard and one risk review.
\item Propose two ADRs for consequential alternatives, without treating an ADR as a viewpoint or a whole architecture description.
\item Defend two conscious diagram omissions.
\end{itemize}

\section{Mastery rubric}
\begin{tabularx}{\linewidth}{@{}YYY@{}}\toprule
\textbf{Dimension} & \textbf{Weight} & \textbf{Evidence} \\ \midrule
Semantic correctness & 30\% & Abstraction, boundary, relationship precision \\
Evidence and traceability & 25\% & Code, runtime and cross-view references \\
Communication & 20\% & Audience clarity and independent interpretation \\
Security and review & 15\% & Threat surface, reviewer findings, resolved blockers \\
Tool reproducibility & 10\% & Versioned source, documented render, lint run \\
\bottomrule\end{tabularx}


\section{Source register}
\begin{itemize}
\item PRIMARY: Simon Brown, The C4 Model: Visualizing Software Architecture, O’Reilly Media, 2026, Chapters 1–12 and Appendix (user-supplied study copy).
\item AUTHORITATIVE C4: https://c4model.com/ — model abstractions, seven diagram types, official review checklist.
\item STANDARD SCOPE: ISO/IEC/IEEE 42010:2022, official catalogue https://www.iso.org/standard/74393.html; conformance requires the full standard.
\item UML: Object Management Group, UML 2.5.1, https://www.omg.org/spec/UML/2.5.1.
\item METHOD: Clements et al., Documenting Software Architectures: Views and Beyond, 2nd ed., Addison-Wesley, 2010/2011.
\item TOOLING: https://plantuml-stdlib.github.io/C4-PlantUML/ and repository tooling/plantuml/README.md.
\end{itemize}

\end{document}
